\documentclass{article}

\usepackage[utf8]{inputenc}
\usepackage[T2A]{fontenc}
\usepackage{helvet}
\usepackage{geometry}
\usepackage{xcolor}
\usepackage{eso-pic}
\usepackage{fancyhdr}
\usepackage{parskip}
\usepackage{microtype}
\usepackage[english, ukrainian]{babel}
\usepackage{url}
\usepackage{amsmath}
\usepackage{amssymb}
\usepackage{amsthm}
\usepackage{longtable}
\usepackage{booktabs}
\usepackage{hyperref}
\usepackage{titlesec}

\definecolor{nato}{HTML}{293757}
\definecolor{footergray}{gray}{0.65}

\geometry{a4paper,left=3cm,right=3cm,top=3cm,bottom=3cm}
\renewcommand{\familydefault}{\sfdefault}
\setcounter{secnumdepth}{3}

\titleformat{\section}
  {\color{nato}\Huge\bfseries}{\thesection.}{1em}{}
\titlespacing*{\section}{0pt}{30pt}{12pt}

\titleformat{\subsection}
  {\color{nato}\LARGE\bfseries}{\thesubsection.}{1em}{}
\titlespacing*{\subsection}{0pt}{20pt}{8pt}

\titleformat{\subsubsection}
  {\color{nato}\Large\bfseries}{\thesubsubsection.}{1em}{}
\titlespacing*{\subsubsection}{0pt}{10pt}{6pt}

\fancyhf{}
\fancyhead[R]{\small\color{footergray} 23 червня 2026}
\fancyfoot[L]{\small\color{footergray} Посібник користувача підсистеми відеоконференцзв’язку ЄСІТС}
\fancyfoot[R]{\small\color{footergray}\thepage}

\newcommand\HUGE[1]{\fontsize{40}{40}\selectfont #1}

\renewcommand{\headrulewidth}{0pt}
\renewcommand{\footrulewidth}{0.4pt}
\renewcommand{\footrule}{\color{footergray}\hrule width \textwidth height 0.4pt}

\begin{document}

\pagestyle{fancy}

\begin{titlepage}
  \AddToShipoutPictureBG*{\AtPageLowerLeft{\color{nato}\rule{\paperwidth}{\paperheight}}}
  \color{white}\thispagestyle{empty}\vspace*{4cm}\centering
  {\Huge  \textbf{ПОСІБНИК КОРИСТУВАЧА} \par} \vspace{2cm}
  {\HUGE  \textbf{Підсистема відеоконференцзв’язку} \par} \vspace{2.5cm}
  {\LARGE \textbf{Єдина судова інформаційно-комунікаційна система} \par} \vspace{2.5cm}
  \vfill
  {\large 2026 \copyright\ ДП «Інформаційні судові системи» \par}
\end{titlepage}

\newpage

\begin{center}
  {\Huge\color{nato}\bfseries Зміст\par}
\end{center}
\vspace{40pt}
\tableofcontents

\newpage

\section{Автори}

Максим Сохацький і Дмитро Лісовий.

\section{Скорочення та визначення}

У цьому документі наведені наступні визначення та скорочення:

\begin{longtable}{p{4cm}p{10cm}}
\toprule
\textbf{Термін / Скорочення} & \textbf{Значення та опис в контексті системи} \\
\midrule
\textbf{ВКЗ} & Підсистема відеоконференцзв'язку (комплекс технічних засобів та програмного забезпечення за посиланням на веб-порталі \url{https://vkz.court.gov.ua}). \\
\textbf{ЄСІТС} & Єдина судова інформаційно-телекоммунікаційна система. \\
\textbf{КЕП} & Кваліфікований електронний підпис (X.509). \\
\textbf{Автентифікація} & Процедура підтвердження електронної ідентифікації користувача за допомогою КЕП в Електронному кабінеті. \\
\textbf{Організатор} & Роль користувача ВКЗ ЄСІТС, що дозволяє створювати відеоконференції (зазвичай співробітник суду). \\
\textbf{Відеофонограма} & Відео- та аудіозапис, утворений під час проведення відеоконференції, завірений КЕП відповідальної особи суду. \\
\textbf{Носій відеозапису} & Оптичний диск або дисковий масив, на який здійснюється запис відеофонограми, який є додатком до протоколу судового засідання, або дисковий масив, на якому здійснюється збереження відеофонограм. \\
\textbf{Lite конференція} & Спрощений режим конференції, який працює за принципом Peer-to-Peer (P2P) без центрального медіа-сервера. \\
\textbf{ЗЗС} & Зал судового засідання, зареєстрований як профіль типу «Приміщення». \\
\bottomrule
\end{longtable}

\section{Загальні положення}

Підсистема відеоконференцзв'язку (ВКЗ) забезпечує учасникам справи можливість брати участь у судовому засіданні в режимі відеоконференції поза межами приміщення суду за допомогою власного Електронного кабінету та власних технічних засобів; у приміщенні іншого суду (за допомогою технічних засобів суду); або в установах виконання покарань, попереднього ув'язнення чи медичних закладах.

ВКЗ також надає можливість громадянам записатися на відеоприйом до обраної організації (наприклад, Державної судової адміністрації України або територіального управління, суду) у зручний час та подати документи на розгляд посадової особи.

Ризики технічної неможливості участі у відеоконференції поза межами приміщення суду (переривання зв’язку тощо) несе учасник справи, який подав відповідну заяву. У разі участі в приміщенні суду — технічну можливість та якість забезпечує відповідний суд.

\newpage

\section{Технічні та мережеві вимоги}

Для забезпечення стабільної роботи ВКЗ необхідно дотримуватися наступних технічних і мережевих параметрів. Блокувальники реклами, VPN-сервіси, проксі-сервери та антивірусні програми можуть перешкоджати нормальній роботі сервісу WebRTC, тому їх необхідно вимкнути або налаштувати перед початком конференції.

\subsection{Підтримувані пристрої та браузери}
\begin{itemize}
    \item \textbf{Комп'ютери (Windows, Linux, macOS):}
    \begin{itemize}
        \item Google Chrome (поточна версія) --- \textbf{рекомендований браузер};
        \item Opera, Mozilla Firefox, Microsoft Edge (поточні версії);
        \item Safari (для macOS).
    \end{itemize}
    \item \textbf{Мобільні пристрої (Android, iOS):} Google Chrome для Android та Safari для iOS.
\end{itemize}

\subsection{Мережеві налаштування (для закритих мереж)}
Для проходження медіа-трафіку WebRTC у корпоративних або закритих мережах на брандмауері мають бути відкритими такі порти:
\begin{itemize}
    \item \textbf{Сигнальний трафік та веб-інтерфейс:}
    \begin{itemize}
        \item TCP-порти \texttt{80, 443, 8443, 8444, 8958} до серверів \url{vkz.court.gov.ua};
        \item TCP-порт \texttt{8958} до сигнального сервера \url{signal-vkz.court.gov.ua}.
    \end{itemize}
    \item \textbf{Медіа-трафік (WebRTC / TURN ноди):}
    \begin{itemize}
        \item TCP/UDP-порти в діапазоні \texttt{31001-52000} до пулу TURN-серверів: \url{turn.court.gov.ua}, \url{turn1.court.gov.ua} --- \url{turn10.court.gov.ua};
        \item TCP-порти \texttt{8443, 8444} до зазначеного пулу TURN-серверів.
    \end{itemize}
    \item \textbf{Мінімальна швидкість Інтернет-з'єднання:} не менше \textbf{2 Mbps} (для обох напрямків).
\end{itemize}

Для спрощених \textbf{Lite конференцій} (працюють по P2P):
\begin{itemize}
    \item TCP-порти \texttt{80, 443, 8958}, TCP/UDP-порти \texttt{3478, 5349};
    \item UDP-порти \texttt{49152-65535} (для прохідного трафіку між клієнтами);
    \item Швидкість Інтернет-з'єднання: не менше \textbf{1 Mbps}.
\end{itemize}

\subsection{Вимоги до обладнання залу судового засідання}
\begin{longtable}{@{}p{3.5cm}p{5.0cm}p{5.5cm}@{}}
\toprule
\textbf{Параметр} & \textbf{Мінімальні вимоги (ED)} & \textbf{Оптимальні вимоги (Full HD)} \\
\midrule
\textbf{Процесор} & Intel Core 2 Duo (2.2 GHz), AMD Athlon 64 X2 або аналог & Intel Core i5 / i7 Sandy Bridge або новіший (від 3.0 GHz) \\
\textbf{ОЗП} & 2 GB & 4 GB або більше \\
\textbf{Відеокарта} & Сумісна з DirectX 9c (256 MB) & Сумісна з DirectX 10.0 (512 MB+) \\
\textbf{Камера} & Logitech C270, HD Pro C910 або аналог & Logitech C920, PTZ-камери HD, плати BlackMagic, AverMedia \\
\bottomrule
\end{longtable}

\newpage

\section{Реєстрація та авторизація користувачів}

\subsection{Для фізичних осіб (профіль «Користувач»)}
Первинна реєстрація та авторизація здійснюється автоматично після переходу до пункту «Відеозв’язок» в Електронному кабінеті Електронного суду. Дані переносяться автоматично. Логіном є адреса електронної пошти користувача.

\subsection{Для співробітників суду та органів правосуддя}
Для реєстрації обов'язково використовується пошта в судовому домені (\url{@court.gov.ua}, \url{@arbitr.gov.ua} тощо). Авторизація також здійснюється через Електронний кабінет Електронного суду.

\subsection{Для залів судових засідань та установ (профіль «Приміщення»)}
\begin{enumerate}
    \item Користувач суду з роллю «Адміністратор організації» створює профіль з типом «Приміщення» у вкладці «Користувачі» службового кабінету.
    \item Зазначається назва або номер залу та судова установа.
    \item На пошту надходить лист для встановлення пароля (мінімум 8 символів: цифри, великі та малі латинські літери, спецсимволи).
    \item Для авторизації залів судових засідань використовується кнопка «Вхід для залів судових засідань, установ...» на сайті ВКЗ.
\end{enumerate}

\subsection{Аварійний режим реєстрації та авторизації}
У разі технічної недоступності авторизації через Електронний суд, вмикається аварійний вхід безпосередньо на порталі ВКЗ за допомогою сервісу електронної ідентифікації \url{Id.gov.ua} (через зчитування КЕП та підтвердження пошти кодом).

\newpage

\section{Організація та проведення конференцій}

Відповідальний співробітник суду повинен авторизуватися на сайті за 10--15 хвилин до початку слухання та перевірити обладнання.

\subsection{Конференція з відеозаписом та веденням протоколу}
\begin{enumerate}
    \item У розділі «Мої конференції» натиснути «Створити конференцію».
    \item Додати учасників конференції за їхніми e-mail або ПІБ.
    \item Увімкнути перемикач «Судове засідання».
    \item У вкладці «Протокол» внести необхідні дані про судову справу.
    \item Натиснути кнопку «Розпочати конференцію».
    \item Виконати обов'язковий тестовий запис тривалістю не менше 10 секунд, перевірити якість звуку та відео.
    \item Під час засідання вносити дії до протоколу.
    \item Після закінчення слухання натиснути кнопку «Завершити конференцію».
\end{enumerate}

\subsection{Конференція з відеозаписом без ведення протоколу}
Дії аналогічні, але перемикач «Судове засідання» не вмикається (або вкладка «Протокол» залишається порожньою). Ведеться лише запис відеофонограми.

\subsection{Конференція без учасників (фіксація засідання в залі суду)}
Застосовується у випадках, коли всі учасники присутні в залі суду особисто. При початку конференції без додавання віддалених користувачів необхідно натиснути «ТАК» у спливаючому запиті «Ви дійсно бажаєте розпочати конференцію без учасників?». Це оптимізує медіа-сервер та виключає ресурсоємне мікшування відеопотоків.

\subsection{Режим аудіофіксації}
Якщо судове засідання фіксується лише в режимі аудіо, після початку конференції організатор вимикає власну камеру на панелі керування.

\subsection{Спрощена Lite конференція}
Призначена для не судових засідань, нарад або робочих зустрічей. Створюється через пункт меню «Створити Lite конференцію». Працює за технологією Peer-to-Peer без запису на серверах системи.

\newpage

\section{Керування учасниками та модерація}

Організатор конференції має повний контроль над ходом відеосесії через панель керування.

\subsection{Панель учасників та модерація}
Кожен учасник відображається у списку з поточним статусом підключення. Організатор може здійснювати такі дії:
\begin{itemize}
    \item \textbf{Вивести на трибуну / Прибрати з трибуни:} Тільки учасники на трибуні транслюють свої відео- та аудіо-потоки іншим користувачам та потрапляють до підсумкового відеозапису.
    \item \textbf{Вимкнути мікрофон / камеру:} Організатор може примусово вимкнути звук або відео будь-якому учаснику для запобігання перешкодам.
    \item \textbf{Заблокувати учасника:} Видалення користувача з кімнати конференції без права повторного входу.
\end{itemize}

\subsection{Налаштування чутливості мікрофона}
У разі використання декількох мікрофонів у залі суду рекомендується увімкнути прапорець «Автоматичне встановлення рівня чутливості мікрофону» в меню налаштувань звуку. Це дозволяє уникнути відлуння (echo-cancelling) та вирівняти гучність промовців.

\section{Голосування та додаткові функції}

\subsection{Проведення голосування під час судового засідання}
\begin{enumerate}
    \item Організатор натискає кнопку «Голосування» на панелі інструментів.
    \item Вводить тему (запитання) та варіанти відповідей (за замовчуванням: «За», «Проти», «Утримався»).
    \item Натискає «Запустити голосування». У всіх учасників на екрані з’являється діалогове вікно.
    \item Після завершення голосування результати автоматично відображаються організатору та можуть бути імпортовані в текстовий протокол засідання.
\end{enumerate}

\subsection{Демонстрація екрана}
Учасники на трибуні мають можливість продемонструвати свій екран або вікно окремого додатка іншим учасникам (наприклад, для пред'явлення електронних доказів чи документів) за допомогою кнопки «Поділитися екраном».

\newpage

\section{Відеозапис та робота з архівами}

\subsection{Запис відеофонограми}
Відеозапис розпочинається автоматично після натискання кнопки «Розпочати конференцію» (за умови увімкненого перемикача запису). Запис ведеться в хмарному сховищі медіа-сервера.

\subsection{Підписання протоколу та блокування змін}
Після завершення конференції секретар судового засідання перевіряє протокол, вносит остаточні правки та підписує його КЕП. Після накладання електронного підпису протокол та пов’язана з ним відеофонограма блокуються від будь-яких подальших змін та автоматично переносяться до центрального сховища Scality.

\subsection{Скачування та онлайн перегляд архіву}
Учасники справи, які сплатили судовий збір, отримують доступ до архіву засідання в Електронному кабінеті. Система підтримує:
\begin{itemize}
    \item \textbf{Пряме скачування файлів:} Завантаження архіву у форматі MP4 та PDF-протоколу.
    \item \textbf{Онлайн перегляд з часовими мітками:} Перегляд відеозапису судового засідання безпосередньо в браузері з можливістю швидкого переходу до конкретної секунди відеофонограми при натисканні на відповідну дію в протоколі засідання.
\end{itemize}

\end{document}
